% ECONOMICS_PROBLEM: {"id": "D44-34", "title": "Optimal auctions beyond independent private values", "jel_code": "D44", "source_id": "OP-0007"}
% !TeX program = xelatex
\documentclass[11pt,a4paper]{article}
\usepackage[margin=23mm]{geometry}
\usepackage{fontspec}
\setmainfont{Latin Modern Roman}
\usepackage{amsmath,amssymb,mathtools}
\usepackage{xcolor,enumitem,booktabs,longtable,array,xurl,hyperref}
\definecolor{muted}{HTML}{53636C}
\hypersetup{colorlinks=true,urlcolor=blue,linkcolor=blue}
\setlength{\parindent}{0pt}
\setlength{\parskip}{5pt}
\newcommand{\R}{\mathbb R}
\newcommand{\E}{\mathbb E}
\newcommand{\N}{\mathbb N}
\newcommand{\ind}{\mathbf 1}
\newcommand{\Var}{\operatorname{Var}}
\newcommand{\Cov}{\operatorname{Cov}}
\newcommand{\supp}{\operatorname{supp}}
\DeclareMathOperator*{\argmin}{arg\,min}
\DeclareMathOperator*{\argmax}{arg\,max}
\newcommand{\entrymeta}[3]{{\sffamily\small\color{muted}\textbf{\texttt{#1}}\quad|\quad #2\par #3\par}}
\newcommand{\Problemref}[1]{\texttt{#1}}
\newcommand{\JELref}[2]{\texttt{#2} (JEL \texttt{#1})}
\begin{document}
\section*{Optimal auctions beyond independent private values}
\textbf{Problem ID:} \texttt{D44-34}\quad \textbf{JEL:} \texttt{D44}\par
% BEGIN ECONOMICS STATEMENT
\label{problem:OP-0007}
\label{atlas:prob:M7}

\noindent\textbf{Original atlas ID:} M7. \textbf{Formulation:} editorial specification of a research family.

\subsection*{Definitions and assumptions}

Fix $n$ bidders, a finite set $K$ of indivisible objects, a feasible allocation set $X$, and finite private signals $s_i\in S_i$. An environment $\omega$ specifies a joint signal law $\pi_\omega$, possibly interdependent values $v_{i,\omega}(x,s)$, nondecreasing utility over money $U_{i,\omega}$, private liquidity bounds $b_{i,\omega}(s_i)$, and entry costs. A mechanism $m$ is a finite extensive-form auction with perfect recall, including any entry and information-acquisition stages, feasible allocation lotteries, and transfer rules. Realized payments must meet the specified budget constraints. Let $\mathcal E(m,\omega)$ be its nonempty set of sequential-equilibrium assessments, using the finite-game consistency definition; every admissible auction and environment must meet this requirement. Admissible mechanisms must satisfy the declared ex-ante or interim participation requirement, with utility evaluated under the environment's actual risk preferences.

\subsection*{Formal research question}

For a stated collection $\Omega$ of such environments and feasible mechanism class $\mathfrak M_F$, characterize
\[
\sup_{m\in\mathfrak M_F}\inf_{\omega\in\Omega}\inf_{e\in\mathcal E(m,\omega)}
 \mathbb E_{m,\omega,e}\!\left[\sum_{i=1}^n t_i-C(x)\right],
\]
where $t_i$ is the bidder-to-seller payment and $C(x)$ the declared seller cost. Investigate restrictions giving a tractable optimal or approximately optimal format when correlation, interdependence, liquidity, risk, and entry coexist. If the objective is efficiency instead, replace revenue with total surplus and explicitly decide how entry and risk-bearing costs enter it. Prove an approximation bound relative to the same feasible class, not to an unrelated independent-private-values optimum.

\subsection*{Interpretation and scope}

This formulation records the source atlas's extensions without pretending that one scalar virtual-value formula covers all of them. Interdependent values mean that utility depends on others' true signals even after allocation; correlation of private values concerns their joint distribution and is a distinct feature. Risk aversion concerns monetary utility and must not be conflated with a hard liquidity constraint. Entry is an equilibrium decision whose response can change both participation and information. A finite-type, quasilinear direct-mechanism problem with fixed entry is frequently expressible as a linear program; that computational benchmark is not being proposed as an unsolved theorem. Continuous types, richer information acquisition, and practical format restrictions require their own guarantees.

\subsection*{Source audit and status}

Myerson (1981), Milgrom--Weber (1982), and Maskin--Riley (2000) are verified through publisher or author primary records. Original link [13] correctly identifies Myerson. Milgrom--Weber analyzes affiliated/interdependent auction environments under restrictions, and Maskin--Riley studies asymmetry; neither provides a universal best auction across all listed features. The shorthand title ``beyond independent private values'' is therefore too broad to define a single theorem-level open problem. The robust extensive-form objective above is an editorial research-family specification. Many subcases have established solutions, and any claimed new open instance must name its primitive restrictions, solution concept, and feasible auction formats.

\subsection*{References}

\begin{enumerate}

\item[S1] Roger B. Myerson (1981). \emph{Optimal Auction Design}. Mathematics of Operations Research 6(1), 58--73. \url{https://pubsonline.informs.org/doi/10.1287/moor.6.1.58}. \textit{Locator:} abstract; optimal-auction characterization. \textit{Establishes:} Revenue-optimal auctions under the paper's private-information assumptions. \textit{Does not establish:} one universally optimal format with all the atlas extensions.

\item[S2] Paul R. Milgrom and Robert J. Weber (1982). \emph{A Theory of Auctions and Competitive Bidding}. Econometrica 50(5), 1089--1122. \url{https://web.stanford.edu/~milgrom/publishedarticles/articlesmain.htm}. \textit{Locator:} publication list and article; affiliated-values analysis. \textit{Establishes:} Auction comparison with interdependence and affiliation assumptions. \textit{Does not establish:} optimal mechanism design for every interdependent-value environment.

\item[S3] Eric Maskin and John Riley (2000). \emph{Asymmetric Auctions}. Review of Economic Studies 67(3), 413--438. \url{https://academic.oup.com/restud/article-abstract/67/3/413/1547420}. \textit{Locator:} abstract; asymmetric bidding distributions. \textit{Establishes:} Revenue comparisons and bidding analysis under asymmetry. \textit{Does not establish:} a universal ranking of auction formats.

\end{enumerate}

\subsection*{Common conventions: Source mathematical conventions}

Unless an entry states otherwise, agent and firm sets are finite. State and action spaces are measurable, strategies and outcome maps are measurable, probability laws are specified, and expectations exist for the targets under discussion. Infinite-horizon discount factors lie in $(0,1)$ and discounted objectives are finite. Norms, tolerances and complexity input sizes must be specified for computational comparisons. Constants or primitives that appear locally are inputs, not empirical facts asserted by this catalogue.

\textbf{Identification.} A statistical model $m\in\mathcal M$ implies an observable law $P_m$ and target $\theta(m)$. At law $P$, its identified set is
\[
\Theta(P;\mathcal M)=\{\theta(m):m\in\mathcal M,\ P_m=P\}.
\]
Point identification holds when this set is a singleton, or equivalently when $P_{m_1}=P_{m_2}$ implies $\theta(m_1)=\theta(m_2)$ for all compatible models. Sharp bounds mean bounds attained, or approached when the boundary is not attained, by compatible models. Writing this set is a definition, not a solution: the substantive task is to establish informative restrictions and derive or estimate the set. An unrestricted-confounding model may give only trivial bounds.

\textbf{Inference.} Identification, estimability and valid uncertainty quantification are separate questions. Fix $\alpha\in(0,1)$. A confidence procedure $C_n$ over a declared law class $\mathcal P$ has asymptotic uniform coverage at level $1-\alpha$ when
\[
\liminf_{n\to\infty}\inf_{P\in\mathcal P}\Pr_P\{\theta(P)\in C_n\}\geq1-\alpha.
\]
For set-identified targets the coverage event must be defined separately, for example coverage of the full identified set. If an entry uses identification-indexed models instead of uniquely indexed laws, the same coverage convention is applied over those models. Pointwise limits do not establish uniform validity, and asymptotic validity does not guarantee finite-sample precision. Sample sizes, dependence structures and weak-identification sequences are part of each application.

\textbf{Counterfactuals.} $Y_i(a)$ is a potential outcome under a specified intervention, including its timing, implementation and versions. It is not $Y_i$ conditional on voluntarily choosing $a$. A full intervention vector permits interference; shorthand scalar treatment notation does not silently impose no interference. Assignment, selection, attrition, anticipation and measurement must be specified. A claim about an instrument's local effect is not automatically a population policy effect.

\textbf{Economic equilibrium.} A counterfactual model includes best responses, constraints, feasibility, market clearing, state transitions and expectations consistent with the stated information. When several equilibria exist, either a declared selector is an input or the target is set-valued. Comparisons across policy regimes must use the stated selection convention. Reduced-form relationships are not presumed invariant after an equilibrium-changing intervention.

\textbf{Optimization and welfare.} Optimization is over the stated feasible set. If attainment is not established, $\sup$ or $\inf$ and $\varepsilon$-optimal choices replace an asserted maximizer or minimizer. For example, with value $V=\sup_{a\in\mathcal A}W(a)<\infty$, an $\varepsilon$-optimal set is $\{a\in\mathcal A:W(a)\geq V-\varepsilon\}$ for $\varepsilon>0$. Benefits, costs, risk and health or environmental outcomes can be aggregated only using declared units and conversion weights. Interpersonal utility weights, the treatment of fiscal transfers and foreign profits, discounting, and the behavioral welfare criterion are explicit inputs. A comparative-static derivative additionally requires existence and appropriate differentiability of the selected counterfactual.

\textbf{Sources.} Local labels [S1], [S2], and so forth restart in every question. Locators identify the relevant abstract, model, section or result at the level actually checked; a bibliographic or abstract-level verification is not represented as inspection of an entire proof. Working papers are distinguished from later publications and changing versions. Source summaries are paraphrased. All URL checks and editorial formulations refer to the audit date stated above.
% END ECONOMICS STATEMENT
\end{document}
